\section{MATERIALES}

Para la construcción del edificio se utilizará hormigón armado en los elementos estructurales principales, tales como muros, vigas y losas. La resistencia del hormigón se ha definido según la demanda estructural de los niveles:

\begin{itemize}
    \item Subterráneo y niveles inferiores: Se utilizará un hormigón de resistencia $f'_c = 35$ MPa.
    \item Pisos superiores: Se utilizará un hormigón de resistencia $f'_c = 25$ MPa.
\end{itemize}

Respecto al acero de refuerzo, se especifica el uso de acero grado A630-420H. Esto implica que el material presenta las siguientes propiedades mecánicas:
\begin{itemize}
    \item Límite de fluencia ($f_y$): $420$ MPa.
    \item Resistencia máxima a la tracción ($f_u$): $630$ MPa.
\end{itemize}

\subsection{Tabiquería y Metalcom}
Para las divisiones interiores y muros no estructurales, se empleará el sistema de perfiles de acero galvanizado (Metalcom), el cual presenta las siguientes características:
\begin{itemize}
    \item Estructura: Perfiles de acero galvanizado conformados en frío.
    \item Revestimiento: Planchas de yeso-cartón (Volcanita) de 10 mm o 15 mm según requerimiento acústico.
    \item Aislación: Lana de vidrio o poliestireno expandido en el interior del tabique para confort térmico.
\end{itemize}

\subsection{Madera}
Se utilizará madera principalmente en elementos de terminación, estructuras de techumbre menores o elementos decorativos, bajo las siguientes especificaciones:
\begin{itemize}
    \item Tipo: Madera de Pino Radiata, clasificada visual o mecánicamente.
    \item Contenido de Humedad: Máximo $19\%$ al momento de la instalación (Madera Seca).
    \item Tratamiento: Aplicación de preservantes contra termitas y humedad, según la norma NCh789.
\end{itemize}